\documentclass{article}
\usepackage{amsmath}
\newtheorem{theorem}{Theorem}
\newtheorem{lemma}[theorem]{Lemma}
\newtheorem{corollary}{Corollary}

\title{Doubling Properties of the Natural Numbers}
\author{A. Author}
\date{2026}

\begin{document}
\maketitle

\begin{abstract}
We study the map $d : \mathbb{N} \to \mathbb{N}$, $d(n) = 2n$, and prove that it
preserves order and stays even. These elementary facts illustrate the
RepoProver paper pipeline on a two-page source.
\end{abstract}

\section{Introduction}\label{sec:intro}

Throughout this paper, $n$ ranges over the natural numbers $\mathbb{N}$. The
central object is the doubling map
\[
d(n) = 2n.
\]
The main result of this paper is Theorem~\ref{thm:double-add}, proved in
Section~\ref{sec:results}. We also use the classical fact that multiplication
is monotone \cite{lean-mathlib}.

\begin{lemma}\label{lem:double-zero}
$d(0) = 0$.
\end{lemma}

\begin{proof}
$d(0) = 2 \cdot 0 = 0$.
\end{proof}

\begin{theorem}[Doubling is strictly monotone]\label{thm:double-strict}
For all $n \in \mathbb{N}$, $n < d(n)$.
\end{theorem}

\begin{proof}
$n < 2n$ because $0 < n$.
\end{proof}


\section{Results}\label{sec:results}

\begin{theorem}[Doubling is additive]\label{thm:double-add}
For all $n, m \in \mathbb{N}$, $d(n + m) = d(n) + d(m)$.
\end{theorem}

\begin{proof}
$d(n + m) = 2(n + m) = 2n + 2m = d(n) + d(m)$ by distributivity.
\end{proof}

\begin{corollary}\label{cor:double-two}
$d(2) = 4$.
\end{corollary}

\begin{lemma}\label{lem:double-even}
For all $n \in \mathbb{N}$, $d(n)$ is even.
\end{lemma}

\begin{proof}
$d(n) = 2n = 2 \cdot n$ is even by definition.
\end{proof}

\begin{corollary}
Every $d(n)$ with $n > 0$ is at least $2$.
\end{corollary}

\appendix

\section{Background on Order}\label{app:order}

\begin{lemma}\label{lem:double-mono}
If $n \le m$ then $d(n) \le d(m)$.
\end{lemma}

\begin{proof}
Multiplication by $2$ is monotone.
\end{proof}

\bibliographystyle{plain}
\bibliography{refs}

\end{document}
